% \section{Learning Algorithms and Numerical Verification}
\subsection{Neural Architecture and Training System}
\label{sec:architecture}

The neural components shared by both QLBS and RLOP embody two design
principles that go beyond a straightforward application of actor-critic
methods.
First, rather than regressing a target quantity directly, both networks embed
a known closed-form function as a differentiable layer inside the
computational graph, letting gradient information flow end-to-end through
that layer.
Second, RLOP introduces a dedicated third network whose training signal
derives from full-trajectory error, not per-step advantage, yielding a
genuinely three-component learning system.

\paragraph{Differentiable Analytical Layer.}
A na\"ive design would train the value network to regress option prices
directly.
This is numerically problematic: prices span several orders of magnitude
across strike and maturity, the regression target is regime-sensitive, and
the network receives no guidance about the geometric structure of the
pricing surface.

Instead, each network is designed to output a single positive scalar
$\hat\sigma$ --- a \emph{latent volatility} --- which is then passed into
a closed-form pricing formula $\mathcal{P}$ implemented as a differentiable
PyTorch module:
\begin{equation}
  \hat{V}(x;\,\theta)
  = \mathcal{P}\!\bigl(S,\,K,\,\tau,\,r,\;f_\theta(x)\bigr),
  \label{eq:latent-vol}
\end{equation}
where $x$ is the input state and $f_\theta:\mathbb{R}^9\to\mathbb{R}_{>0}$
is the residual network.
Because $\mathcal{P}$ is differentiable in all its arguments, gradients
flow end-to-end from pricing errors back into $\theta$ without any
approximation.
This design is directly analogous to \emph{physics-informed neural networks}
(PINNs) \cite{raissi2019physics}: rather than learning the entire
input--output map from scratch, the network learns only the correction
that the analytical model cannot capture, while the closed-form handles the
well-understood structure.
As a result, the effective network output is interpretable --- it recovers
the implied-volatility surface as an emergent by-product of training,
not as a post-hoc diagnostic.

\paragraph{CombinedResNet with Embedded Input Normalisation.}
The nine-dimensional raw state
$(S_t,\, t,\, T\!-\!t,\, K,\, r,\, \mu,\, \sigma,\, \lambda,\, c)$
contains quantities whose natural scale varies by orders of magnitude across
training instances.
Rather than handling this as a separate preprocessing step, we embed the
following coordinate transforms directly into the network's \texttt{forward}
pass, ensuring they are always applied consistently during both training
and inference:
\begin{equation}
  \tilde{m} = \frac{\log(S_t/K) + r\,(T\!-\!t)}{\sigma\sqrt{T\!-\!t}},
  \quad
  \tilde{\tau} = \sqrt{\frac{T\!-\!t}{T_{\mathrm{ref}}}},
  \label{eq:input-transform}
\end{equation}
where $\tilde{m}$ is the log-moneyness normalised by the at-the-money
time-value scale, and $T_{\mathrm{ref}}$ is a fixed reference maturity.
Embedding these transforms in the architecture plays the same role as
tokenisation or positional encoding in sequence models: structural knowledge
about the input domain is encoded once at the architecture level rather than
learned from data.

The backbone is a \emph{CombinedResNet} of $D$ stacked residual blocks
\cite{he2016deep} with configurable depth $D\in\{2,3,4\}$.
Each block applies a two-layer MLP with layer normalisation and a skip
connection.
The policy head outputs a Gaussian distribution
$\pi(a\mid x) = \mathcal{N}\!\left(\mu_\pi(x;\theta_\mu),\,
\sigma_\pi(x;\theta_\sigma)^2\right)$,
where $\mu_\pi$ is produced by the CombinedResNet and $\sigma_\pi$ by a
separate branch with a sigmoid clamp,
$\sigma_\pi = 0.01 + 0.90\cdot\mathrm{sigmoid}(g_\theta(x))$,
bounding exploration noise away from both zero and saturation.
A value network of identical architecture provides a learned baseline for
variance reduction.
Both networks are trained with REINFORCE \cite{williams1992simple} and
optimised by Adam at learning rate $10^{-4}$.

\paragraph{RLOP: Three-Network Decoupled Design.}
Standard actor-critic algorithms maintain two learned components: a policy
$\pi_\theta$ and a value baseline $V_\phi$.
RLOP extends this to three networks by introducing a dedicated
\emph{PriceEstimator} $P_\psi$, whose training objective and signal
source differ fundamentally from those of $V_\phi$:
\begin{itemize}[leftmargin=1.2em,topsep=1pt,itemsep=0pt]
  \item \textbf{Policy} $\pi_\theta$: updated by policy gradient using
    per-step advantages computed from $V_\phi$.
  \item \textbf{Value baseline} $V_\phi$: trained on per-step discounted
    returns $G_t$ to reduce gradient variance (standard REINFORCE baseline).
  \item \textbf{PriceEstimator} $P_\psi$: trained once per episode on the
    \emph{full-trajectory replication error}
    $\delta^{(i)} = h(S_T) - \Pi_T^{(i)}$, aggregated across all
    concurrently managed maturities $i=1,\dots,T$.
\end{itemize}
The key insight is that $V_\phi$ and $P_\psi$ answer different questions:
$V_\phi$ estimates the expected cost-to-go from a state, while $P_\psi$
learns the fair replication price from terminal outcomes.
Conflating these two objectives into a single critic, as standard
actor-critic would, degrades training stability when the reward is dominated
by terminal rather than per-step feedback.
The PriceEstimator shares its backbone architecture with the policy,
so gradient information from trajectory-level pricing errors flows
indirectly into $\pi_\theta$ through the shared feature representation,
without requiring an explicit auxiliary loss on the policy.

\paragraph{Scalable Parallel Training and Two-Mode Inference.}
Both environments are implemented as fully vectorised PyTorch modules
supporting $N_{\mathrm{par}}=2{,}500$ parallel simulation paths per
episode.
Each call to \texttt{env.step()} advances all paths simultaneously via
batched tensor operations, eliminating per-episode Python overhead and
enabling GPU-accelerated training.

For downstream evaluation --- computing hedge ratios and option values
across large Monte Carlo ensembles --- the trained networks are loaded in a
\emph{grad-fn inference mode}: gradient computation is disabled and
operations are pinned to CPU, decoupling the memory budget for training from
the throughput requirements of path-wise simulation.
This two-mode design (GPU training / CPU inference) is a practical
engineering choice that keeps wall-clock evaluation time tractable when
sweeping across many contract configurations.

The full ablation covers depth $D\in\{2,3,4\}$, risk-aversion coefficients
$\lambda\in\{0.01,0.1,1.0\}$, and ten independent replicas per
configuration, yielding 90 trained models per method and providing
the statistical confidence intervals reported in \cref{sec:empirical}.
